\section{Introduction}

Reinforcement Learning from Human Feedback (RLHF) is a key technique for aligning Large Language Models~(LLMs) with human values and preferences ~\citep{vemprala2023chatgpt,achiam2023gpt, Ouyang2022, shen2024policy,shen2025exploring,hu2024openrlhf}. Despite the emergence of non-RL alternatives like DPO ~\citep{rafailov2023direct}, state-of-the-art applications such as ChatGPT/GPT-4 ~\citep{vemprala2023chatgpt, openai2023gpt4}, Claude ~\citep{anthropic2023introducing}, and Gemini ~\citep{team2023gemini} continue to rely on RL algorithms.

The dominant RL algorithm in this space is Proximal Policy Optimization (PPO) ~\citep{schulman2017proximal}. PPO employs an "Actor-Critic" architecture, where a dedicated critic network is trained to estimate the advantage function (i.e., the extent to which an action yields a higher expected return than the mean). However, this critic network introduces substantial computational overhead and memory requirements, making training prohibitively expensive and limiting large-scale model alignment on small-scale clusters \citep{shao2024deepseekmath}.

To address PPO's efficiency issues, a family of REINFORCE-based methods has emerged without the critic model, including ReMax~\citep{li2023remax}, RLOO~\citep{2024Back}, and GRPO~\citep{shao2024deepseekmath}. These algorithms replace the critic network with an estimator that computes the advantage from multiple responses to the same prompt.

However, this critic-free approach introduces a critical challenge: \textit{advantage estimation}. GRPO uses \textbf{prompt-level (local) normalization}, estimating a baseline and standard deviation from responses to one prompt; RLOO uses a leave-one-out baseline without inherently requiring this scale normalization. These are distinct design choices. First, Appendix~\ref{appendix:proof_GRPO} shows that finite-sample self-normalization differs from population standardization. Second, small groups (e.g., $k=4$ or $k=8$) make the scale sensitive to sampled outcomes. Homogeneous groups provide no centered reward signal, while different reward spreads induce different prompt-specific scales. Finally, these relative comparisons may emphasize fitting individual groups rather than transfer across prompts. Our experiments examine the resulting training dynamics and held-out generalization.

We introduce two distinct variants of this framework, each tailored to a specific use case. \ours is an efficient algorithm with $k \ge 1$. In its $k=1$ configuration (detailed in Algorithm 1), it maximizes prompt diversity for general-domain RLHF. It can also be applied with $k>1$ (Section 4.2), using the same global normalization principle. \varn is a specialized \textit{variant} for complex tasks ($k>1$) that benefits from \textbf{group sampling}. It fixes the flaws of GRPO by first subtracting the \textit{group mean} (for reward reshaping) and then normalizing by the \textit{global standard deviation} (for stability). It also employs a $k_2$ gradient surrogate for separate $\mathrm{KL}$ regularization, as detailed in Appendix~\ref{appendix:kl_design}. In summary, our contributions are as follows:
\begin{itemize}
    \item We provide a theoretical proof that finite-sample local normalization differs from population standardization under explicit sampling assumptions (Appendix~\ref{appendix:proof_GRPO}).
    \item We propose \ours, a critic-free framework centered on Global Advantage Normalization, as a stable, efficient, and theoretically sound alternative.
    % --- MODIFIED: Clarified the two algorithms based on the new k>=1 logic ---
    \item We present two specialized algorithms: We can use \ours when $k \ge 1$ for efficient general RLHF and reasoning, and use \varn when $k>1$ for robust complex agentic tasks.
    % --- END OF MODIFICATION ---
    \item The experiments show that \ours achieves competitive general RLHF scores with shorter responses and stronger held-out reasoning performance. At the same time, \varn outperforms both GRPO and PPO in complex agentic reasoning tasks.
\end{itemize}